\input{../../common/slide-common-header.tex}
\input{../../common/slide-plan.tex}

\title[]{Lecture \syncPrimitivesNum: \syncPrimitivesTopic}
\subtitle[]{\syncPrimitivesKey}
\author[]{Alexander Filatov\\ filatovaur@gmail.com}

\date{}

\begin{document}

\begin{frame}
  \titlepage
  \url{https://github.com/Svazars/parallel-programming/blob/main/slides/pdf/l4.pdf}
\end{frame}

\begin{frame}{In previous episodes}

\begin{itemize}
    \item We study communication and coordination of threads in pre-emptive multitasking OS
    \item Execution speed is not under our control
    \item Decoupling tasks -- \texttt{Runnable}, \texttt{Callable}, \texttt{Future}
        \begin{itemize}
            \item from particular threads  -- \texttt{ThreadFactory}
            \item from submission protocol -- \texttt{Executor}
            \item from scheduling policy   -- \texttt{ExecutorService}
        \end{itemize}        
    \item Tools for thread coordination:
        \begin{itemize}
            \item \texttt{Thread.join}, \texttt{Lock}, \texttt{Condition}, \texttt{ExponentialBackoff}, \texttt{Future.get}
        \end{itemize}
    \item Concurrency pitfalls
        \begin{itemize}
            \item wait-for graph, deadlock, thread-starvation deadlock, resource deadlock
            \item fairness, priority inversion, starvation
            \item lost signal, spurious wakeup, predicate invalidation
        \end{itemize}    
\end{itemize}
\end{frame}

\begin{frame}{Lecture plan}
\tableofcontents
\end{frame}

%\begin{frame}[t]{Disclaimer}
%
%Concurrent programming is evolving branch of Computer Science.
%
%\pause
%
%There are many designs for key synchronization primitives:
%
%\begin{itemize}
    %\item ''Monitor classification'' by Buhr et al\footnote<2->{\tiny\url{https://dl.acm.org/doi/10.1145/214037.214100}}
    %\item ''Semaphores in Plan 9'' by Mullender and Cox\footnote<2->{\tiny\url{https://swtch.com/semaphore.pdf}}
    %\item ''Futexes are tricky'' by Drepper\footnote<2->{\tiny\url{https://www.semanticscholar.org/paper/Futexes-Are-Tricky-Drepper/%86253463161f45ea013a501f89c8de36f1a09192}}
%\end{itemize}
%
%\pause
%We will stick to the Java programming language (built-in monitors) and \texttt{java.util.concurrent} package.
%
%\pause
%It is enough to illustrate design space, trade-offs, common correctness problems.
%
%\pause
% It is \textbf{not} enough to start concurrency-oriented low-level development in other languages (e.g. C++ stdlib/boost, Golang %with plan9 style sema etc).
%
%\end{frame}

\section{Monitor}

\subsection{API}

\begin{frame}[t,fragile]{Mutex: popular use-cases}

\begin{minted}{java}
interface Lock {
    void lock();
    void unlock();
}
\end{minted}
\end{frame}

\questiontime{OOP, modularity, code reusage -- which mutex property helps to support it?}

\begin{frame}[t,fragile,noframenumbering]{Mutex: popular use-cases}

\begin{minted}{java}
interface Lock {
    void lock();
    void unlock();
}
\end{minted}

\begin{itemize}
    \item Reentrant
\end{itemize}

\end{frame}

\questiontime{why do we develop multithreaded software?}

\begin{frame}[t,fragile,noframenumbering]{Mutex: popular use-cases}

\begin{minted}{java}
interface Lock {
    void lock();
    void unlock();
}
\end{minted}

\begin{itemize}
    \item Reentrant
    \item Throughput-oriented
\end{itemize}

\end{frame}

\questiontime{fair or unfair mutex helps to achieve better throughput of system-as-a-whole?}

\begin{frame}[t,fragile,noframenumbering]{Mutex: popular use-cases}

\begin{minted}{java}
interface Lock {
    void lock();
    void unlock();
}
\end{minted}

\begin{itemize}
    \item Reentrant
    \item Unfair
\end{itemize}
\end{frame}

\questiontime{mutex solves problem of thread \textbf{contention}. What should we use to simplify thread \textbf{coordination}?}

\begin{frame}[t,fragile,noframenumbering]{Mutex: popular use-cases}

\begin{minted}{java}
interface Lock {
    void lock();
    void unlock();
    Condition newCondition();
}
\end{minted}

\begin{itemize}
    \item Reentrant
    \item Unfair
\end{itemize}

\begin{minted}{java}
interface Condition {
    void await();
    void signal();
    void signalAll();
}
\end{minted}
\end{frame}

\questiontime{how many Condition Variables do you need for the most use cases we studied so far?}

\begin{frame}[t,fragile,noframenumbering]{Mutex: popular use-cases}

\begin{minted}{java}
interface Lock {
    void lock();
    void unlock();
    Condition condition();
}
\end{minted}

\begin{itemize}
    \item Reentrant
    \item Unfair
    \item Single condition
\end{itemize}

\begin{minted}{java}
interface Condition {
    void await();
    void signal();
    void signalAll();
}
\end{minted}
\end{frame}

\begin{frame}[t,fragile,noframenumbering]{Mutex: popular use-cases}

\begin{minted}{java}
interface PopularLock {
    void lock();
    void unlock();

    void await();
    void signal();
    void signalAll();
}
\end{minted}

\begin{itemize}
    \item Reentrant
    \item Unfair
    \item Single condition
\end{itemize}
\end{frame}

\questiontime{control flow and mutex -- any pitfalls?}

\begin{frame}[t,fragile,noframenumbering]{Mutex: popular use-cases}

\begin{minted}{java}
interface PopularLock {
    void lock();
    void unlock();

    void await();
    void signal();
    void signalAll();
}
\end{minted}

\begin{itemize}
    \item Reentrant
    \item Unfair
    \item Single condition
    \item Structured locking: \texttt{lock try-finally unlock}
\end{itemize}
\end{frame}

\questiontime{how to strictily enforce structured locking on source code level?}

\begin{frame}[t,fragile,noframenumbering]{Mutex: popular use-cases}

\begin{minted}{java}
interface PopularLock {
    // lock-unlock provided by language keyword
    void await();
    void signal();
    void signalAll();
}
\end{minted}

\begin{itemize}
    \item Reentrant
    \item Unfair
    \item Single condition
    \item Structured locking: language support
\end{itemize}
\end{frame}


\begin{frame}[t,fragile,noframenumbering]{Mutex: popular use-cases}

\begin{minted}{java}
interface PopularLock {
    // lock-unlock provided by language keyword
    void await();     // <-- this
    void signal();    //          could
    void signalAll(); //                clash
}
\end{minted}

\begin{itemize}
    \item Reentrant
    \item Unfair
    \item Single condition
    \item Structured locking: language support
\end{itemize}
\end{frame}

\begin{frame}[t,fragile,noframenumbering]{Mutex: popular use-cases}
\begin{minted}{java}
interface PopularLock {
    // lock-unlock provided by language keyword
    void wait();     
    void notify();   
    void notifyAll();
}
\end{minted}

\begin{itemize}
    \item Reentrant
    \item Unfair
    \item Single condition
    \item Structured locking: language support
\end{itemize}
\end{frame}

\showTOCSub

\begin{frame}[t,fragile]{Monitor}

\begin{minted}{java}
class Monitor { void wait();
                void notify();
                void notifyAll();
}
\end{minted}

\begin{itemize}
    \item Reentrant, Unfair
    \item \texttt{synchronized} for critical section (always structured)
\end{itemize}

\begin{minted}{java}
void main() {
    var monitor = ...;
    synchronized (monitor) {
        foo();
        if (bar()) { monitor.wait(); } 
        else       { monitor.notify(); }    
    }
}
\end{minted}
\end{frame}

\questiontime{in which Java package \texttt{Monitor} class is declared?}

\begin{frame}[t,fragile]{Built-in Java monitor}

\begin{minted}{java}
java.lang.Object {
    void wait();
    void notify();
    void notifyAll();
}
\end{minted}

\begin{itemize}
    \item Reentrant, Unfair
    \item \texttt{synchronized} for critical section (always structured)
    \item \textbf{Every} Java object has associated monitor
\end{itemize}

{\small\url{https://docs.oracle.com/en/java/javase/11/docs/api/java.base/java/lang/Object.html}}

\pause
It \textbf{does not} mean you allocate monitor for every object, though
\begin{itemize}
    \item ''Thin locks: featherweight synchronization for Java'' by Bacon et al. 1998\footnote<2->{\tiny\url{https://dl.acm.org/doi/abs/10.1145/277652.277734}}
\end{itemize}
\end{frame}


\begin{frame}[t,fragile]{Disclaimer}

Concurrent programming is rapidly evolving branch of Computer Science.

\pause
There are many designs for key synchronization primitives:

\begin{itemize}
    \item ''Monitor classification'' by Buhr et al\footnote<2->{\tiny\url{https://dl.acm.org/doi/10.1145/214037.214100}}
    \item ''Semaphores in Plan 9'' by Mullender and Cox\footnote<2->{\tiny\url{https://swtch.com/semaphore.pdf}}
    \item ''Futexes are tricky'' by Drepper\footnote<2->{\tiny\url{https://dept-info.labri.fr/~denis/Enseignement/2008-IR/Articles/01-futex.pdf}}
\end{itemize}

\pause
Today we will stick to the Java programming language (built-in monitors) and \texttt{java.util.concurrent} package.

\pause
It is enough to illustrate design space, trade-offs, common correctness problems.

\end{frame}


\begin{frame}[fragile]{Java monitor}
\framesubtitle{Design}

Effectively \texttt{ReentrantLock + single Condition}

\begin{itemize}
    \item \texttt{MonitorEnter}\footnote{\tiny\url{https://docs.oracle.com/javase/specs/jvms/se11/html/jvms-6.html\#jvms-6.5.monitorenter}}\textsuperscript{,}\footnote{\tiny\url{https://docs.oracle.com/javase/specs/jvms/se11/html/jvms-3.html\#jvms-3.14}}
    \item \texttt{MonitorExit}
    \item \texttt{wait}
    \item \texttt{notify}
    \item \texttt{notifyAll}
\end{itemize}

\begin{tikzpicture}[remember picture,overlay]
\node[xshift=3.5cm,yshift=-2.5cm] at (current page.center) {\includegraphics[width=0.27\textwidth]{./pics/MonitorJava.png}};
\end{tikzpicture}

\end{frame}


\begin{frame}[fragile]{Toy problem: CompletionService}

{\tiny\url{https://docs.oracle.com/en/java/javase/11/docs/api/java.base/java/util/concurrent/CompletionService.html}}
\begin{minted}{java}
interface CompletionService<V> {
    Future<V> submit(Callable<V> task);
    Future<V> take();
}
\end{minted}
\begin{itemize}
    \item \texttt{take} retrieves the \texttt{Future} representing the next completed task, blocks if no available
    \item thread-safe dequeue
    \begin{itemize}
        \item producers submit tasks
        \item consumers process completed tasks
    \end{itemize}
\end{itemize}
\end{frame}

\begin{frame}[fragile,noframenumbering]{CompletionService via built-in monitors}
\begin{minted}{java}
class MyCompletionService<V> implements CompletionService<V> {
  ExecutorService service; Queue<V> q = new LinkedList<>();
  MyCompletionService(ExecutorService service) { this.service = service; }    
  Future<V> submit(Callable<V> task) {
    return service.submit(() -> {
      V result = task.call();
      synchronized (q) { q.add(result); q.notify(); }
    });}
  Future<V> take() {
    synchronized (q) { 
      while (q.isEmpty()) { q.wait(); }
      return new AlreadyCompletedFuture<>(q.remove());
}}}
\end{minted}
\end{frame}


\begin{frame}[fragile]{\texttt{Monitor}ing at home}

{\small\url{https://github.com/Svazars/parallel-programming/tree/main/hw/block1/4.1}}
\begin{homeworkmail}{Task 4.1}
    \begin{itemize}
        \item Find bugs in small programs that use built-in monitors
    \end{itemize}
\end{homeworkmail}

\end{frame}

\begin{frame}[fragile]{Java monitor: for every object}

\textbf{Every} Java object have associated built-in monitor.

\begin{minted}{java}
  Object x = new Object();
  synchronized(x) {
    x.wait();
    x.notifyAll();
  }
\end{minted}

\begin{itemize}
    \item Special keyword: \texttt{synchronized} (no \texttt{.lock/unlock}, no try-finally, always structured)
    \item \texttt{Condition} and \texttt{ReentrantLock} are Java objects and they have built-in monitor
    \item Library-level synchronization and language-level synchronization are \textbf{independent}, do not mix \texttt{notify/signal} and \texttt{wait/await}.
\end{itemize}

\end{frame}

\begin{frame}[t,fragile]{Java monitor: spot-a-problem №1}

\begin{minted}{java}
  Lock l = ...;
  synchronized (l) {
    foo();
  }
\end{minted}

\pause 

\texttt{synchronized} -> \texttt{monitorenter/monitorexit} 

\pause

\texttt{monitorenter != Lock.lock}

\end{frame}

\begin{frame}[t,fragile]{Java monitor: spot-a-problem №2}

\begin{minted}{java}
  Lock l = ..;
  l.lock();
  try {
    while (!bar()) { 
      l.wait(); 
    }
    baz();
  } finally { 
    l.unlock(); 
  }
\end{minted}

\pause 

\texttt{wait} -> built-in monitor associated with \texttt{l}, not some \texttt{Condition}

\end{frame}

\begin{frame}[fragile]{Java monitor: spot-a-problem №3}

\begin{minted}{java}
  Lock l = ..;
  Condition c = l.newCondition();
  ...
  l.lock();
  try {
    while (!bar()) { 
      c.wait(); 
    }
    baz();
  } finally { 
    l.unlock(); 
  }
\end{minted}

\pause
\texttt{wait} -> built-in monitor associated with \texttt{c}, not \texttt{Condition.await}

\end{frame}


\begin{frame}[t,fragile]{Java monitor: for every method}

\textbf{Every} Java method could be marked as \texttt{synchronized}.
\begin{minted}{java}
  public void foo() { doStuff(); }
  public synchronized void bar() { doStuffUnderMutex(); }
\end{minted}

It synchronizes on \texttt{this} (instance methods) or \texttt{Class<?>} (static methods):
\begin{minted}{java}
  public synchronized void fooSignature() { doStuff(); }
  public void fooInternal() {
    synchronized(this) {
      doStuff();
    }
  }
\end{minted}
\end{frame}

\begin{frame}[t,fragile]{Java monitor: spot-a-problem №4}

\begin{minted}{java}
class MyClass {
  private final Object guard = new Object();
  public void foo() { synchronized(guard) { bar(); } }
  public synchronized void bar() { }
}
\end{minted}

\end{frame}

\begin{frame}[t,fragile,noframenumbering]{Java monitor: spot-a-problem №4}

\begin{minted}{java}
class MyClass {
  private final Object guard = new Object();
  public void foo() { synchronized(guard) { bar(); } }
  public synchronized void bar() { }
}
class OhMyClass extends MyClass {
  public synchronized void baz() { 
    foo();
  }
}
\end{minted}

\pause
Possible deadlock:
\begin{itemize}
    \item Thread A calls \texttt{foo}: \texttt{lock(guard) -> lock(this)}
    \item Thread B calls \textt{baz}: \texttt{lock(this) -> lock(guard)}
\end{itemize}
\end{frame}

\questiontime{Is it good or bad to mark methods as thread-safe (\texttt{synchronized} on publicly available signature level)?}

\begin{frame}[t,fragile]{Java monitor: for every method}

\textbf{Every} Java method could be marked as \texttt{synchronized}.
\begin{minted}{java}
  public void foo() { doStuff(); }
  public synchronized void bar() { doStuffunderMutex(); }
\end{minted}

Inheritance or \texttt{synchronized(myAwesomeClassInstance)}
\begin{itemize}
    \item may violate your locking policy 
    \item may cause a deadlock    
\end{itemize}

It is common pattern to use \texttt{private final Object lock} and do not expose it as public API.
\end{frame}


\begin{frame}[t,noframenumbering]{Java monitor}
\framesubtitle{Common pitfalls}

Correctness/safety:
\begin{itemize}
    \item Blocking \texttt{MonitorEnter}, \texttt{wait} -- deadlock, locking order
    \item Reentrant
    \item Visibility/consistency -- \texttt{synchronized} looks like ''atomic transaction''
    \item Lost signal
    \item Predicate invalidation
    \item Spurious wakeup
\end{itemize}

Fairness/starvation:
\begin{itemize}
    \item Admission policy for \texttt{MonitorEnter}/\texttt{MonitorExit} is not specified
    \item Admission policy for \texttt{wait}/\texttt{notify} is not specified
\end{itemize}
\end{frame}

\questiontime{what is spurious wakeup for built-in Java monitor?}

\begin{frame}[t]{Homebrewed with love}

{\small\url{https://github.com/Svazars/parallel-programming/tree/main/hw/block1/4.2}}
\begin{homeworkmail}{Task 4.2}
    \begin{itemize}
        \item Implement \texttt{Monitor} using \texttt{ReentrantLock} and \texttt{Condition}
        \item \textbf{Hard}: implement it without spurious wakeups
    \end{itemize}
\end{homeworkmail}
\end{frame}


\begin{frame}[t,fragile]{Java monitor}
\framesubtitle{Specific pitfalls}

Problem: unintended mix of \texttt{java.util.concurrent} and built-in synchronization
\begin{itemize}
    \item \texttt{wait} vs \texttt{await}
    \item \texttt{synchronized} vs \texttt{lock/unlock}
\end{itemize}
Solution: static analysis (FindBugs/InteliJ IDEA)

\pause
\ 
\newline

Problem: unintended encapsulation breach
\begin{itemize}
    \item \texttt{synchronized} methods inheritance
    \item \texttt{synchronized} on object outside of library method
\end{itemize}

Solution: \texttt{private final} guards

\end{frame}

\begin{frame}{Java monitor}
\framesubtitle{Rules of thumb}

Mutual exclusion
\begin{itemize}
    \item If you need ''just mutex'' -- consider using encapsulated object + \texttt{synchronized}
    \begin{itemize}
        \item Save your time with try-finally
        \item Other stuff (built-ins have better monitoring and optimizer-friendly) is debatable
    \end{itemize}

    \pause

    \item Use \texttt{Lock} if you need ''advanced'' properties:
    \begin{itemize}
        \item Fairness (remember, scheduling is counter-intuitive)
        \item Non-structured locking (remember, exceptions are hard)
        \item Polling via \texttt{tryLock} (remember, race conditions are tricky)
    \end{itemize}
\end{itemize}

\pause
Signalling\pause. The story is absolutely the same.

Built-in monitors are ''friendlier'' yet less flexible.
\end{frame}


\begin{frame}[t,fragile]{Signalling: broadcast}

\begin{itemize}
    \item \texttt{Lock+Condition} and \texttt{Monitor} provide \texttt{notifyAll}
\end{itemize}

\begin{tikzpicture}[remember picture,overlay]
\node[xshift=3.5cm,yshift=-6.5cm] at (current page.center) {\includegraphics[width=0.27\textwidth]{./pics/MonitorJava.png}};
\end{tikzpicture}

\end{frame}

\questiontime{Does \texttt{notifyAll} or \texttt{signalAll} actually force all waiting threads to wake-up?}

\begin{frame}[t,noframenumbering]{Signalling: broadcast}

\begin{itemize}
    \item \texttt{Lock+Condition} and \texttt{Monitor} provide \texttt{notifyAll}
    \item Effectively threads are ''released'' one-by-one, since they are forced to enter critical section.    
    \pause \item Seems that encourages excessive context switch
\end{itemize}
\end{frame}

\begin{frame}[fragile]{Thundering herd problem}

    Scenario:
    \begin{itemize}
        \item many threads wake-up because of broadcast 
        \item only few (e.g. single one) make progress
    \end{itemize}

    \pause
    Problem: inefficient utilization of scheduling quantums

    \pause

    Solution:
    \begin{itemize}
        \item Use broadcasting with care
        \pause \item \texttt{monitor.notifyAll} and \texttt{condition.signalAll} could avoid this by waking threads one-by-one?
    \end{itemize}
\end{frame}

\subsection{Wait Morphing}
\showTOCSub

\begin{frame}[t,fragile]{Disclaimer}

\begin{itemize}
    \item Code on the next slides is a \textbf{sketch} implementation of \texttt{Monitor} class
    \item It excessively uses mutexes and polling with timeout to make the code more compact
    \item It does not support \texttt{interrupt}s as built-in Java monitors (see Lecture~\patternsNum)
    \item Real world is much more complicated
    \begin{itemize}
        \item ''Compact Java Monitors'' by Dice, Kogan. 2021\footnote{\tiny\url{https://arxiv.org/pdf/2102.04188}}
    \end{itemize}
\end{itemize}

We just want to understand how ''awake other thread'' operation could be delayed until mutex is unlocked.

\end{frame}


\begin{frame}[t,fragile]{Wait morphing}

\begin{minted}{java}
class Monitor {    
  Thread owner; 
  int count;
  Set<Thread> enterSet; 
  static Thead me() { return Thread.currentThread(); }
  synchronized boolean tryEnter() {
    if (owner == me()) { count++; return true; }
    if (owner == null) { count = 1; owner = me(); return true; }
    return false;
  }
  void monitorEnter() { ... }    
  void monitorExit()  { ... }  
}
\end{minted}
\end{frame}

\begin{frame}[t,fragile,noframenumbering]{Wait morphing}

\begin{minted}{java}
class Monitor {  
  void monitorEnter() {
    synchronized (this) { enterSet.add(me()); }
    while (!tryEnter()) { suspendMe(timeout=100ms); }
  }
  void monitorExit() { Thread e = null;
    synchronized(this) { 
      if (owner != me()) { throw new IllegalMonitorStateException(); }
      count--;
      if (count == 0) { enterSet.remove(me()); owner = null;
                        e = selectRandom(enterSet);
      }
    }
    if (e != null) { awake(e); }
  }
}
\end{minted}
\end{frame}

\begin{frame}[t,fragile,noframenumbering]{Wait morphing}

\begin{minted}{java}
class Monitor { Set<Thread> waitSet; Set<Thread enterSet; Thread owner;
  void wait() {
    if (owner != me()) { throw new IllegalMonitorStateException(); }
    waitSet.add(me());    
    while (waitSet.contains(me())) { 
      int cnt = fullyUnlock();
      suspendMe(timeout=100ms); 
      relock(cnt);
  }}
  void notify() {
    if (owner != me()) { throw new IllegalMonitorStateException(); }
    Thread w = removeRandom(waitSet);
    if (w != null) { awake(w); }
  }
}
\end{minted}
\end{frame}

\begin{frame}[t,fragile,noframenumbering]{Wait morphing}

\begin{minted}{java}
class Monitor { Set<Thread> waitSet; Set<Thread enterSet; Thread owner;
  void wait() {
    if (owner != me()) { throw new IllegalMonitorStateException(); }
    waitSet.add(me());    
    while (waitSet.contains(me())) { 
      int cnt = fullyUnlock();
      suspendMe(timeout=100ms); 
      relock(cnt);
  }}
  void notify() {
    if (owner != me()) { throw new IllegalMonitorStateException(); }
    Thread w = removeRandom(waitSet); // vvvvvvvvvvvvvvvvv
    if (w != null) { synchronized(this) { enterSet.add(w); } }
  }
}
\end{minted}
\end{frame}

% TODO: summary Thundering Herd problem appears for broadcast signalling with subsequent serialization bottleneck (e.g. mutex)
%       javas monitor.notifyAll/condition.signalAll avoid this problem by using wait morphing
%             they know the bottleneck place (mutex to reacquire)
%             they avoid context switch by delaying thread awakening (awake only one thread on unlock)
%

\questiontime{wait morphing is cool! But what if my application really needs to awaken many threads simultaneously?}

\section{Broadcast}
\subsection{CountDownLatch}
\showTOCSub

\begin{frame}[fragile]{CountDownLatch}
    \begin{center} \includegraphics[width=0.8\textwidth]{./pics/horse_race.png} \end{center}
\end{frame}

\begin{frame}[fragile,noframenumbering]{CountDownLatch}

CountDownLatch\footnote{\tiny\url{https://docs.oracle.com/en/java/javase/11/docs/api/java.base/java/util/concurrent/CountDownLatch.html}}

\begin{minted}{java}
  CountDownLatch(int count)
  void await()     // awaits counter == 0
  void countDown() // counter = max(counter - 1, 0)
  long getCount()  // return counter
\end{minted}

Common usage patterns:
\begin{itemize}
    \item \texttt{CountDownLatch(1)} acts as a ''gate'':
    \begin{itemize}
      \item \textbf{closed} (\texttt{counter == 1}) 
      \item \textbf{open} (\texttt{counter == 0})
     \end{itemize}

    \item \texttt{CountDownLatch(N)} acts as a ''all threads ready'' point:
    \begin{itemize}
      \item \textbf{not-yet-ready} (\texttt{counter > 0})
      \item \textbf{ready} (\texttt{counter == 0})
     \end{itemize}
\end{itemize}

\end{frame}


\begin{frame}[fragile]{CountDownLatch}
\framesubtitle{Example}

\begin{minted}{java}
static CountDownLatch ready = new CountDownLatch(N + 1);
static CountDownLatch start = new CountDownLatch(1);
void main() {
  for (i = 0; i < N; i++) new Thread(new Worker(i)).start();
  ready.countDown(); ready.await();
  start.countDown();
}
void worker() {
  prepare();
  ready.countDown();
  start.await();
}
\end{minted}
\end{frame}

\begin{frame}[fragile]{NaiveCountDownLatch}

\begin{minted}{java}
class NaiveCountDownLatch {
  private int count;
  public CountDownLatch(int count) { 
    if (count < 0) throw new IllegalArgumentException();
    this.count = count; 
  }
  synchronized void await()     { while (count != 0) { this.wait(); } }
  synchronized long getCount()  { return count; }
  synchronized void countDown() {    
    if (count == 1) { count = 0; this.notifyAll(); } 
    else            { count--; }
  }
\end{minted}
\end{frame}


\begin{frame}[fragile]{Homework: efficient CountDownLatch}

{\small\url{https://github.com/Svazars/parallel-programming/tree/main/hw/block1/4.3}}
\begin{homeworkmail}{Task 4.3}
    \begin{itemize}
        \item Implement \texttt{CountDownLatch} using \texttt{ReentrantLock}/\texttt{Condition}/built-in monitors
        \item Ensure it releases threads simultaneously (avoid wait morphing side-effects)
        \item Use condition-per-waiting-thread pattern
    \end{itemize}
\end{homeworkmail}
\end{frame}

% 
% \subsection{CyclicBarrier}
% \showTOCSub
% 
% \begin{frame}[t,fragile]{CyclicBarrier}
% 
% \texttt{CountDownLatch} is single-use because counter monotonically decreases. 
% 
% Some tasks require \texttt{N} iterations of some parallel activity before completion.
% \end{frame}
% 
% \questiontime{Provide examples of parallelizable tasks that consist of \texttt{N} repeating steps.}
% 
% 
% \begin{frame}[t,fragile,noframenumbering]{CyclicBarrier}
% 
% \texttt{CountDownLatch} is single-use because counter monotonically decreases. 
% 
% Some tasks require \texttt{N} iterations of some parallel activity before completion.
% 
% Creating \texttt{ArrayList<CountDownLatch>} is cumbersome.
% 
% \end{frame}
% 
% 
% \begin{frame}[fragile]{CyclicBarrier}
% 
% CyclicBarrier\footnote{\tiny\url{https://docs.oracle.com/en/java/javase/11/docs/api/java.base/java/util/concurrent/CyclicBarrier.html}}
% 
% \begin{minted}{java}
%   CyclicBarrier(int parties, Runnable barrierAction)
%   int await()     
%   int getNumberWaiting()  
%   int getParties()    
% \end{minted}
% 
% Common usage patterns: mostly-parallel tasks with sequential coordination
% \begin{itemize}
%     \item Multiply matrices row-by-row (parallel), handle result (sequential)
%     \item Do numerical simulation for \texttt{dt} (parallel), save intermediate result (sequential)
%     \item Handle batch of requests (parallel), decide next batch size (sequential)
% \end{itemize}
% 
% % HW: implement via single monitor
% \end{frame}
% 
% \begin{frame}[fragile]{CyclicBarrier}
% \framesubtitle{Example}
% 
% They have example with matrix multiplication right in the javadoc\footnote{{\tiny\url{https://docs.oracle.com/en/java/javase/11/docs/api/java.base/java/util/concurrent/CyclicBarrier.html}}}!
% \end{frame}


\subsection{Thundering herd problem}
\showTOCSub

\begin{frame}[fragile,noframenumbering]{Thundering herd problem}

\begin{minted}{java}
static Deque<Object> data = ...
static CountDownLatch dataReady = new CountDownLatch(1);
void main(Collection<Object> workItems) {
  synchronized(data) {  data.addAll(workItems); }
  dataReady.countDown();
}
void worker() {
  dataReady.await();
  Object element;
  synchronized(data) { element = data.poll(); }
  process(element);
}
\end{minted}
\end{frame}

\begin{frame}[fragile,noframenumbering]{Thundering herd problem}

    Scenario:
    \begin{itemize}
        \item many threads wake-up because of broadcast 
        \item only few (e.g. single one) make progress
    \end{itemize}

    Problem: inefficient utilization of scheduling quantums.

    Solution:
    \begin{itemize}
        \item Use broadcasting with care
    \end{itemize}
\end{frame}

\questiontime{I am confused, could I observe thundering herd problem when using monitors or could not I?}

\section{Group-level concurrency}
\showTOC

\begin{frame}{Group level concurrency}

\begin{itemize}
    \item \texttt{Lock/Monitor} divides threads into ''owner'' and ''others''.
    \item \texttt{Condition} divides threads into ''waiters'' and ''others''.
    \item \texttt{CountDownLatch} helps to coordinate start/stop for group of threads.
\end{itemize}

Could we do more?

\begin{itemize}
    \pause \item No more than N threads in critical section \pause (semaphore)
    \pause \item One exclusive writer and N independent readers \pause (read-write lock)
\end{itemize}
\end{frame}

\subsection{Semaphore}
\showTOCSub

\begin{frame}[fragile]{Semaphore}

Semaphore\footnote{\tiny\url{https://docs.oracle.com/en/java/javase/11/docs/api/java.base/java/util/concurrent/Semaphore.html}}

\begin{minted}{java}
  Semaphore(int permits)
  void acquire() // acquire permit, blocks until available
  void release() // release permit
\end{minted}

Important notes:
\begin{itemize}
    \item Perfect to 
    \begin{itemize}
        \item limit number of active threads (parallelism control) 
        \item bound shared resource usage (concurrent access control)
    \end{itemize}    
    \item No ownership -- any thread could \texttt{release} permits    
\end{itemize}
\end{frame}

\questiontime{\texttt{Semaphore} could limit number of simultaneously executing threads to \texttt{N == 1}. \newline Have we already seen similar behaviour?}

\questiontime{\texttt{Semaphore} could limit number of simultaneously executing threads to \texttt{N > 1}. \newline Have we already seen similar behaviour?}

\begin{frame}[fragile]{BinaryLock}

\begin{minted}{java}
class BinaryLock { // NonReentrantLock
  Semaphore semaphore = new Semaphore(1);
  Thread owner = null;
  void lock() {
    semaphore.acquire();
    owner = Thread.currentThread();
  }
  void unlock() {
    if (owner != Thread.currentThread()) {
      throw new IllegalMonitorStateException();
    }
    owner = null;
    semaphore.release();
  }
}
\end{minted}
\end{frame}

\begin{frame}[fragile]{Resource control}

\begin{minted}{java}
static Semaphore sema = new Semaphore(N);
static Set<Object> items = ...;
void useItem() {
  sema.acquire();
  try {
    synchronized(items) {
      Object item = selectRandom(items);
      items.remove(item);
    }
    use(item);
    synchronized(items) { items.add(item); }
  } finally { sema.release(); }
}
\end{minted}
\end{frame}

\begin{frame}[fragile]{Your own railroad}

{\small\url{https://github.com/Svazars/parallel-programming/tree/main/hw/block1/4.4}}
\begin{homeworkmail}{Task 4.4}
    \begin{itemize}
        \item Implement \texttt{Semaphore} using \texttt{ReentrantLock}/\texttt{Condition}/built-in monitors        
    \end{itemize}
\end{homeworkmail}
\end{frame}

\subsection{ReadWriteLock}
\showTOCSub

\begin{frame}[t,fragile]{ReadWriteLock}

ReadWriteLock\footnote{\tiny\url{https://docs.oracle.com/en/java/javase/11/docs/api/java.base/java/util/concurrent/locks/ReadWriteLock.html}}

\begin{minted}{java}
  Lock readLock();
  Lock writeLock();
\end{minted}

Read-mostly data structures:
\begin{itemize}
    \item Concurrent readers, Exclusive writer
\end{itemize}

Design decisions:
\end{frame}

\questiontime{What are the main design decisions for any mutual exclusion primitive?}

\begin{frame}[t,fragile,noframenumbering]{ReadWriteLock}

ReadWriteLock\footnote{\tiny\url{https://docs.oracle.com/en/java/javase/11/docs/api/java.base/java/util/concurrent/locks/ReadWriteLock.html}}

\begin{minted}{java}
  Lock readLock()
  Lock writeLock()
\end{minted}

Divides threads into two groups:
\begin{itemize}
    \item ''Simultaneous readers'' (many)
    \item ''Exclusive writer'' (only one)
\end{itemize}

Design decisions:
\begin{itemize}
    \item Reentrancy for writers, reentrancy for readers

    \item Reader preference vs. Writer preference
%    TODO-q1: if N active readers why to stop N+1-th reader
%    TODO-q2: if there are X waiting writers and Y waiting readers who should be preferred

    \item Lock upgrade (promotion): readLock -> writeLock

    \item Lock downgrade: writeLock -> readLock    
\end{itemize}
\end{frame}


\begin{frame}[fragile]{ReadWriteLock}
\framesubtitle{Example}

\begin{minted}{java}
static Map<K, V> kv = ...;
static ReadWriteLock rw = ...;
V get(K k) {
  rw.readLock().lock();
  try {
    return kv.get(k);
  } finally { rw.readLock().unlock(); }
}
void put(K k, V v) {
  rw.writeLock().lock();
  try {
    kv.put(k, v);
  } finally { rw.writeLock().unlock(); }
}
\end{minted}
\end{frame}

% \begin{frame}[fragile]{Homework}
% TODO: h/w implement ReadWriteLock using built-in lock/condition/monitors
% \end{frame}

\questiontime{suppose I have method \texttt{foo(int x)} and cannot change the signature. How could I pass additional parameter to the method?}

\questiontime{suppose I have method \texttt{foo(int x)} and cannot change the signature. How could I pass additional parameter to the method \textbf{in concurrent world}?}

\section{ThreadLocal}
\showTOCSub

\subsection{API}

\begin{frame}[fragile]{ThreadLocal}

ThreadLocal\footnote{\tiny\url{https://docs.oracle.com/en/java/javase/11/docs/api/java.base/java/lang/ThreadLocal.html}}
\begin{minted}{java}
class ThreadId { static long nextId = 0; static Object guard = new Object();  
  static final ThreadLocal<Integer> threadId = new ThreadLocal<Integer>() {
    @Override protected Integer initialValue() {
      synchronized(guard) { return nextId++; };
    }
  };
  public static long get() { return threadId.get(); }
}
void main() {
  long myId = ThreadId.get();
  threadFactory.newThread(() -> { long hisId = ThreadId.get(); });
}
\end{minted}
\end{frame}

\subsection{Replication}
\showTOCSub

\begin{frame}[fragile]{Prime numbers: pure function}

\begin{minted}{java}
class PureMath {
  static boolean checkPrime(long num) {
    if (num <  2) { throw new IllegalArgumentException(); }
    if (num <= 3) { return true; }
    for (long t = 2; t < sqrt(num) + 1; t++) {
      if (num % t == 0) { return false; }
    }
    return true;
  }
}
\end{minted}
\end{frame}

\begin{frame}[fragile,noframenumbering]{Prime numbers: caching}

\begin{minted}{java}
class CachingMath {
  static Set<Long> knownPrimes = new Set<>();
  static boolean isPrime(long num) {
    if (knownPrimes.contains(num)) { 
      return true; 
    } else {
      boolean prime = checkPrime(num);
      if (prime) { knownPrimes.add(num); }
      return prime;
    }
  }
}
\end{minted}
\end{frame}

\begin{frame}[fragile,noframenumbering]{Prime numbers: concurrent caching}

\begin{minted}{java}
class ThreadSafeMath {
  static Set<Long> knownPrimes = new Set<>();
  static ReadWriteLock rw = new ReentrantReadWriteLock();
  static boolean isPrime(long num) {
    rw.readLock().lock();
    try     { if (knownPrimes.contains(num)) { return true; } } 
    finally { rw.readLock().unlock(); }
    boolean isPrime = checkPrime(num);
    if (isPrime) {
      rw.writeLock().lock();
      try     { knownPrimes.add(num); } 
      finally { rw.writeLock().unlock(); }
}}}
\end{minted}
\end{frame}

\begin{frame}[fragile,noframenumbering]{Prime numbers: replication}
\begin{minted}{java}
class ReplicatedMath {
  static final ThreadLocal<Set<Long>> myPrimes = new ThreadLocal<>() {
    @Override protected Set<Long> initialValue() { return new HashSet<>(); }
  };
  static boolean isPrime(long num) {
    if (myPrimes.get().contains(num)) { return true; }
    rw.readLock().lock();
    try     { myPrimes.addAll(knownPrimes); } 
    finally { rw.readLock().unlock(); }
    if (myPrimes.get().contains(num)) { return true; }
    boolean prime = checkPrime(num);
    if (prime) { updateKnownPrimes(num); }
    return prime;
}}
\end{minted}
\end{frame}

\begin{frame}[fragile]{Replication}

Idea: 
\begin{itemize}
    \item create thread-local copies of shared data
    \item prefer using ''local cache'' without any synchronization
    \item in-thread operations would be consistent with each other
\end{itemize}

Potential problems:
\begin{itemize}
    \item Cross-thread operations could be partially inconsistent 
    \begin{itemize}
      \item global state ''lags'' (thread-local copy is more informative)
      \item thread-private states ''lag'' (this thread misses info from other threads)
    \end{itemize}    
\end{itemize}

\end{frame}

\questiontime{ \texttt{myPrimes.addAll(knownPrimes)} executed under read lock. Scalable?}

\subsection{Snapshotting}
\showTOCSub

\begin{frame}[t,fragile]{Snapshotting}
\begin{minted}{java}
class PrimeSetSnapshot {
  private final Set<Long> primes;
  public final long version;
  PrimeSetSnapshot(Set<Long> p, long v) {
    this.primes = Collections.unmodifiableSet(new HashSet<>(p));
    this.version = v;
  }
  static PrimeSetSnapshot initial() {
    return new PrimeSetSnapshot(Collections.emptySet(), 0);
  }
  boolean contains(long num) { return primes.contains(num); }  
  PrimeSetSnapshot with(long num) {
    return new PrimeSetSnapshot(new HashSet<>(primes).add(num), version + 1);
  }
}
\end{minted}
\end{frame}

\begin{frame}[t,fragile]{Snapshotting}
\begin{minted}{java}
class SnapshottingMath {
  private PrimeSetSnapshot latest = PrimeSetSnapshot.initial();
  private ReadWriteLock rw = new ReentrantReadWriteLock();  
  static final ThreadLocal<PrimeSetSnapshot> myPrimes = new ThreadLocal<>() {
    @Override protected Set<PrimeSetSnapshot> initialValue() {
      return PrimeSetSnapshot.initial();
  }};
}
\end{minted}

\end{frame}

\begin{frame}[t,fragile,noframenumbering]{Snapshotting}
\begin{minted}{java}
class SnapshottingMath {
  PrimeSetSnapshot latest; ThreadLocal<PrimeSetSnapshot> myPrimes;
  boolean isPrime(long num) {
    ???
}}
\end{minted}
\end{frame}

\begin{frame}[t,fragile,noframenumbering]{Snapshotting}
\begin{minted}{java}
class SnapshottingMath {
  PrimeSetSnapshot latest; ThreadLocal<PrimeSetSnapshot> myPrimes;
  boolean isPrime(long num) {
    PrimeSetSnapshot priv = myPrimes.get(); PrimeSetSnapshot publ = null;
    if (priv.contains(num)) { return true; }    
    rw.readLock().lock(); try { publ = latest; } finally { rw.readLock().unlock(); }
    if (priv.version != publ.version) { 
      myPrimes.set(publ); return isPrime(num);
    } else {
      boolean prime = checkPrime(num);
      if (prime) { updateLatestSnapshot(priv.with(num)); }
      return prime;
    }
}}
\end{minted}
\end{frame}

\begin{frame}[t,fragile]{Snapshotting}

Idea: 
\begin{itemize}
    \item use versioning for shared data: immutable data structure + version number
    \item use immutable snapshot for speculative computation
    \item if your snapshot is up-to-date -- ''commit'' computation result
    \item retry otherwise    
\end{itemize}

Potential problems:
\begin{itemize}
    \item Missing update: shared data changed but version number was not increased
    \item Starvation: multiple retries
    \item Memory consumption: too many heavyweight copies
    \item Internal mutability: speculative computation has side effects
\end{itemize}

Further reading: seqlock, ABA problem
\end{frame}

\begin{frame}[t,fragile]{But they'll never take our FREEDOM! Deadlock-freedom}
% TODO: wallace pic

{\small\url{https://github.com/Svazars/parallel-programming/tree/main/hw/block1/4.5}}
\begin{homeworkmail}{Task 4.5}
    \begin{itemize}
        \item Deadlock empire. Every. Single. Level.
    \end{itemize}
\end{homeworkmail}
\end{frame}


\section{Thread-safe counters}
\showTOCSub

\begin{frame}[t,fragile]{Synchronized counter}

\begin{minted}{java}
final class SynchronizedCounter {
  private long data;
  public SynchronizedCounter(long initial) { 
    data = initial; 
  }
  public synchronized long get() { 
    return data; 
  }
  public synchronized void add(long delta) { 
    data += delta; 
  }
}
\end{minted}
\end{frame}

\begin{frame}[fragile]{Read-mostly counter}

\begin{minted}{java}
final class ReadMostlyCounter {
  private long data;
  private ReadWriteLock rw = new ReentrantReadWriteLock();
  public ReadMostlyCounter(long initial) { data = initial; }
  public long get() { 
    rw.readLock().lock();
    try { return data; } 
    finally { rw.readLock().unlock(); }
  }
  public void add(long delta) { 
    rw.writeLock().lock();
    try { data += delta; } 
    finally { rw.writeLock().unlock(); }
  }
}
\end{minted}
\end{frame}

\begin{frame}[fragile]{Write-mostly counter: lock splitting}

\begin{minted}{java}
final class SynchronizedCounter {
  private long data;
  public SynchronizedCounter(long initial) { 
    data = initial; 
  }
  public synchronized long get() { 
    return data; 
  }
  public synchronized void add(long delta) { // <------ lock splitting
    data += delta; 
  }
}
\end{minted}

\textbf{Reminder}: Homework Task 2.4 (Medium)
\end{frame}

\begin{frame}[t,fragile]{Write-mostly counter: distributed counter}

Distributed counter:
\begin{itemize}
    \item \texttt{N} independent \texttt{long} slots -- 1 slot per thread
    \item Every user of distributed counter should register itself
    \item Thread \texttt{add}s to its own slot (no contention)
    \item Thread \texttt{get}s from all slots (could take time)
\end{itemize}
\end{frame}

\questiontime{''1 slot per thread''. Which Java class should I use?}
\questiontime{''user of distributed counter should register itself''. Should it be reenterable operation?}
\questiontime{
    \begin{itemize}
        \item Thread \texttt{add}s to its own slot (no contention)
        \item Thread \texttt{get}s from all slots (could take time)
    \end{itemize}
    Do you know data structure that divides threads into two groups: ''many simultaneous'' and ''only one''?
}

\begin{frame}[t,fragile,noframenumbering]{Distributed counter: thread-private slots}

\begin{itemize}
    \item \texttt{N} independent \texttt{long} slots -- 1 slot per thread    
\end{itemize}

\begin{minted}{java}
class DistributedCounter {
  class Slot { long data; }
  static final ThreadLocal<Slot> mySlot = new ThreadLocal<>() {
    protected Set<PrimeSetSnapshot> initialValue() { return null; }
  };
  boolean thisThreadIsRegistered() { 
    return mySlot.get() != null; 
  }
  void ensureRegistered() {
    if (!thisThreadIsRegistered()) { mySlot.set(new Slot()); }
}}
\end{minted}

\pause

\begin{itemize}
    \item Thread-safe?
\end{itemize}

\end{frame}


\begin{frame}[t,fragile,noframenumbering]{Distributed counter: thread-specific lazy initialization}

\begin{itemize}
    \item Every user of distributed counter should register itself
\end{itemize}

\begin{minted}{java}
class DistributedCounter {
  ...
  void add(long delta) {
    ensureRegistered();
    ???
    ???
    ???
  }
}
\end{minted}
\end{frame}


\begin{frame}[t,fragile,noframenumbering]{Distributed counter: scalable add}

\begin{itemize}
    \item Thread \texttt{add}s to its own slot (no contention)    
\end{itemize}

\begin{minted}{java}
class DistributedCounter {
  final ReadWriteLock rw = new ReentrantReadWriteLock();
  void add(long delta) {
    ensureRegistered();
    ???
    ???
    ???
  }
}
\end{minted}
\end{frame}

\begin{frame}[t,fragile,noframenumbering]{Distributed counter: scalable add}

\begin{itemize}
    \item Thread \texttt{add}s to its own slot (no contention)    
\end{itemize}

\begin{minted}{java}
class DistributedCounter {
  final ReadWriteLock rw = new ReentrantReadWriteLock();
  void add(long delta) {
    ensureRegistered();
    rw.readLock().lock();
    try { mySlot.get().data++;} 
    finally { rw.readLock().unlock(); }
  }
}
\end{minted}
\end{frame}

\begin{frame}[t,fragile,noframenumbering]{Distributed counter: consistent get}

\begin{itemize}
    \item Thread \texttt{get}s from all slots (could take time)
\end{itemize}

\begin{minted}{java}
class DistributedCounter {
  final ReadWriteLock rw = new ReentrantReadWriteLock();
  long get() {
    ???
    ???
    ???
  }
}
\end{minted}
\end{frame}


\begin{frame}[t,fragile,noframenumbering]{Distributed counter: consistent get}

\begin{itemize}
    \item Thread \texttt{get}s from all slots (could take time)
\end{itemize}

\begin{minted}{java}
class DistributedCounter {
  final ReadWriteLock rw = new ReentrantReadWriteLock();
  long get() { rw.writeLock().lock();
               try { ??? } 
               finally { rw.writeLock().unlock(); }
             }
}
\end{minted}
\end{frame}

\begin{frame}[t,fragile,noframenumbering]{Distributed counter: consistent get}

\begin{itemize}
    \item Thread \texttt{get}s from all slots (could take time)
\end{itemize}

\begin{minted}{java}
class DistributedCounter {
  final ReadWriteLock rw = new ReentrantReadWriteLock();
  long get() { rw.writeLock().lock();
               try { ??? }  
               finally { rw.writeLock().unlock(); }
             }
  final List<Slot> lst = new ArrayList<>();
  void ensureRegistered() {
    if (thisThreadIsRegistered()) { return; }
    mySlot.set(new Slot());
    rw.writeLock().lock(); 
    try { lst.add(mySlot.get()); } finally { rw.writeLock().unlock(); }
}}
\end{minted}
\end{frame}

\begin{frame}[t,fragile,noframenumbering]{Distributed counter: consistent get}

\begin{itemize}    
    \item Thread \texttt{get}s from all slots (could take time)
\end{itemize}

\begin{minted}{java}
class DistributedCounter {
  final ReadWriteLock rw = new ReentrantReadWriteLock();
  long get() { rw.writeLock().lock(); long result = 0;
               try { for (var slot: lst) { result += slot.data; } }  
               finally { rw.writeLock().unlock(); }
               return result; }
  final List<Slot> lst = new ArrayList<>();
  void ensureRegistered() {
    if (thisThreadIsRegistered()) { return; }
    mySlot.set(new Slot());
    rw.writeLock().lock(); 
    try { lst.add(mySlot.get()); } finally { rw.writeLock().unlock(); }
}}
\end{minted}
\end{frame}

\begin{frame}[t,fragile]{Distributed counter: takeaways}

\begin{itemize}
    \item Read-mostly and write-mostly designs could be quite different
    \item \texttt{ThreadLocal} helps to organize uncontended execution
    \item \texttt{ReadWriteLock} splits threads into two groups
    \begin{itemize}
        \item Not necessarily readers and writers!
    \end{itemize}    
\end{itemize}

\pause
Teaser 
\pause
\begin{itemize}
    \item Replication + Snapshotting + Distributed systems = Headache 
    \pause 
    \begin{itemize}
        \item Wait-free snapshots (Lecture~\foundationsPlusNum)
        \item ABA problem (Lecture~\atomicsNum)
        \item Replication on CPU level (Lecture~\cacheCoherencyNum)
    \end{itemize}
\end{itemize}
\end{frame}


% \section{Workload generation}
% \showTOC
% 
% \begin{frame}{Producer-consumer}
% 
% \begin{itemize}
%     \item N threads generate work items
%     \item M threads process work items
% \end{itemize}
% 
% \begin{center}
% \includegraphics[width=0.7\textwidth]{./pics/prod-cons.png}
% \end{center}
% 
% BlockingQueue\footnote{\tiny\url{https://docs.oracle.com/en/java/javase/11/docs/api/java.base/java/util/concurrent/BlockingQueue.html}}
% 
% \end{frame}
% 
% \questiontime{How to ensure that producers and consumers are ''balanced''?}
% 
% 
% \begin{frame}{Producer-consumer}
% 
% \begin{itemize}
%     \item X worker threads
% \end{itemize}
% 
% Depending on heuristic, every thread either
% \begin{itemize}
%     \item generate work item
%     \item process existing work item
% \end{itemize}
% 
% \pause
% 
% Common problems:
% \begin{itemize}
%     \item Fixed-size buffer deadlock: 
%     \begin{itemize} 
%         \item buffer capacity is X
%         \item every thread produces no more than 2 elements per iteration
%     \end{itemize}
% 
%     \item Empty-buffer deadlock: 
%     \begin{itemize} 
%         \item buffer is empty
%         \item every thread heuristically decided to be consumer
%     \end{itemize}
% \end{itemize}
% 
% \end{frame}
% 
% \questiontime{How to ensure that producers and consumers reach ''consistent progress''?}
% 
% \begin{frame}[t,fragile]{Fork-join}
% 
% Fork-join model\footnote{\url{https://en.wikipedia.org/wiki/Fork-join_model}}
% 
% \begin{center}
% \includegraphics[width=0.65\textwidth]{./pics/fork-join.png}
% \end{center}
% 
% \end{frame}
% 
% \questiontime{How to decide when to stop \texttt{fork}?}
% 
% 
% \begin{frame}{Reactive Streams}
% 
% \url{https://openjdk.org/jeps/266}
% 
% \begin{quote}
% The main goal of Reactive Streams is to govern the exchange of stream data across an asynchronous boundary—think passing elements on to another thread or thread-pool
% — while ensuring that the receiving side is not forced to buffer arbitrary amounts of data. In other words, back pressure is an integral part of this model in order to 
% allow the queues which mediate between threads to be bounded. 
% \end{quote}
% 
% \texttt{Flow.Publisher<T>}, \texttt{Flow.Subscriber<T>}\footnote{\tiny\url{https://docs.oracle.com/en/java/javase/11/docs/api/java.base/java/util/concurrent/Flow.html}}
% \end{frame}
% 
% 
% \begin{frame}{Workload generation: takeaways}
% 
% \begin{itemize}
%     \item Producer-consumer: ''push'' based concurrent system
%     \item Fork-join: divide-and-conquer, ''helping hand'' approach
%     \item Reactive: ''propagate backpressure'' design 
% \end{itemize}
% 
% More features -- harder to implement, monitor and control.
% 
% \end{frame}
% 
% 
% \section{Load balancing}
% \showTOC
% 
% 
% \begin{frame}{Load balancing}
% 
% ThreadPool-based design uses classic timesharing approach: 
% \begin{itemize}
%     \item there are \texttt{X} tasks (functions) 
%     \item \texttt{Y} ''virtual'' executors (threads) which are mapped to
%     \item \texttt{Z} ''real'' executors (cores)
% \end{itemize}
% \pause
% There are well-known corner cases:
% \begin{itemize}
%     \item \texttt{X < Y}: insufficient concurrency (coarse-grained problem)
%     \item \texttt{X >> Y}: excessive coordination overheads (too fine-grained approach or too aggressive sub-task creation)
%     \item \texttt{Y < Z}: undersubscription (underutilization)
%     \item \texttt{Y >> Z}: oversubscription
% \end{itemize}
% \pause
% In any scenario, \texttt{Y} threads have to coordinate task execution:
% \begin{itemize}
%     \item Eventual progress: every task will be executed
%     \item Correctness: every task is executed exactly once
%     \item Performance: keep coordination overheads minimal  
% \end{itemize}
% 
% % TODO: arcticle with perf-oriented designs for over/under scubscription
% 
% \end{frame}
% 
% 
% \begin{frame}{Load balancing: global queue}
% \framesubtitle{Work arbitrage}
% 
% Load balancing design:
% \begin{itemize}
%     \item Single shared thread-safe task queue
%     \item Every available \texttt{Worker} gets tasks one-by-one
% \end{itemize}
% 
% Counter-examples:
% \begin{itemize}
%     \item Many short tasks
%     \item \texttt{Worker}s of varying speeds (energy-efficient and usual cores)
% \end{itemize}
% 
% \end{frame}
% 
% \begin{frame}{Load balancing: global queue + batching}
% \framesubtitle{Work arbitrage}
% 
% Load balancing design:
% \begin{itemize}
%     \item Single shared thread-safe task queue
%     \item Every available \texttt{Worker} gets tasks in batches (N per request)
% \end{itemize}
% 
% Counter-examples:
% \begin{itemize}
%     \item Mix of short and long tasks
%     \item \texttt{Worker}s of varying speeds (energy-efficient and usual cores)
% \end{itemize}
% 
% \end{frame}
% 
% 
% \begin{frame}{Load balancing: global queue + batching}
% \framesubtitle{Work dealing}
% 
% Load balancing design:
% \begin{itemize}
%     \item Single shared thread-safe task queue
%     \item Every available \texttt{Worker} gets tasks in batches (N per request)
%     \item Overloaded \texttt{Worker} moves fraction of tasks to global queue
% \end{itemize}
% 
% Counter-examples:
% \begin{itemize}
%     \item Few long tasks
%     \item Overloaded thread spends CPU for auxiliary coordination
% \end{itemize}
% \end{frame}
% 
% \begin{frame}{Load balancing: work-stealing}
% 
% Load balancing design:
% \begin{itemize}
%     \item Several shared thread-safe task queues (e.g. 1 per \texttt{Worker})
%     \item \texttt{Worker} get tasks in batches (N per request) from some queue (''steals'' from neighbour)
% 
% \end{itemize}
% 
% Not that easy to implement and fine-tune\footnote{\tiny\url{https://shipilev.net/#forkjoin}}. Nontrivial termination protocol.
% 
% ForkJoinPool\footnote{\tiny\url{https://docs.oracle.com/en/java/javase/11/docs/api/java.base/java/util/concurrent/ForkJoinPool.html}}
% 
% \end{frame}


\section{Summary}

\begin{frame}{Summary}

Mutual exclusion and signalling combined into single universal primitive:
\begin{itemize}
    \item \texttt{Monitor}
    \item Wait morphing -- scheduler-aware optimization
\end{itemize}

Java language has special support for monitor-based thread coordination:
\begin{itemize}
    \item \texttt{synchronized} block
    \item every object has built-in monitor
    \item any method could be marked as \texttt{synchronized}
\end{itemize}

Bulk thread notification:
\begin{itemize}
    \item \texttt{CountDownLatch}
    \item Thundering herd problem
\end{itemize}

Tools for fine-grained group-level thread control:
\begin{itemize}
    \item \texttt{Semaphore}, \texttt{ReadWritelock}
\end{itemize}

Thread-specific data could be used for caching:
\begin{itemize}
    \item replication, snapshotting
\end{itemize}
\end{frame}


\begin{frame}{Summary: homework}
    \texttt{Lecture 4, Tasks 4.*:} {\tiny\url{https://github.com/Svazars/parallel-programming/blob/main/hw/block1}}
\end{frame}

\end{document}
